% !TEX root = main.tex
\section{Giới thiệu}\label{sec:intro}

\subsection{Bối cảnh và động lực nghiên cứu}\label{sec:intro:context}

\textbf{Lĩnh vực ứng dụng.} Bài toán tìm đường (route planning) là một trong những bài toán kinh điển nhất của khoa học
máy tính và đồng thời là thành phần lõi của các hệ thống giao thông thông minh (Intelligent Transportation Systems):
ứng dụng bản đồ và dẫn đường trên điện thoại, nền tảng gọi xe công nghệ, điều phối xe giao hàng chặng cuối, quản lý
đội xe buýt hay xe cứu thương. Ở dạng đơn giản nhất, mạng đường được mô hình hoá thành một đồ thị có hướng có trọng
số và câu hỏi đặt ra là tìm đường đi có tổng trọng số nhỏ nhất giữa hai đỉnh; thuật toán Dijkstra~\cite{dijkstra1959}
và A*~\cite{hart1968} giải trọn vẹn câu hỏi đó trong thời gian đa thức.

Trong thực tế đô thị, giả thiết ``trọng số cố định'' lại hiếm khi đúng. Thời gian để đi qua một đoạn đường phụ thuộc
mạnh vào \emph{thời điểm} đi vào đoạn đó: cùng một trục đường có thể thông thoáng lúc ba giờ sáng nhưng ùn tắc kéo dài
lúc tám giờ sáng. Hệ quả là tuyến đường ngắn nhất theo chiều dài không nhất thiết là tuyến nhanh nhất, và tuyến nhanh
nhất lúc này không nhất thiết là tuyến nhanh nhất một giờ sau. Đây chính là bài toán \emph{đường đi ngắn nhất phụ thuộc
thời gian} (time-dependent shortest path, TDSP)~\cite{cooke1966,dreyfus1969,delling2009}.

\textbf{Nhu cầu và vấn đề thực tiễn.} Hà Đông là quận cửa ngõ phía tây nam của Hà Nội (trước khi sắp xếp lại đơn vị
hành chính năm 2025), nơi tập trung nhiều khu đô thị mới và các trục hướng tâm chịu lưu lượng lớn như Quang Trung --
Nguyễn Trãi, Tố Hữu (Lê Văn Lương kéo dài), Văn Khê và Quốc lộ~6. Ùn tắc ở khu vực này mang ba đặc điểm điển hình của
giao thông đô thị Việt Nam:
\begin{itemize}
    \item \emph{tính chu kỳ theo ngày}: ùn tắc lặp lại vào các khung giờ cao điểm sáng và chiều;
    \item \emph{tính có hướng}: buổi sáng dòng xe đổ vào trung tâm thành phố, buổi chiều theo chiều ngược lại;
    \item \emph{tính bất định}: ngày này khác ngày kia, và có các sự cố bất thường (va chạm, ngập, sửa đường) làm một
    đoạn đường chậm hẳn trong một khoảng thời gian.
\end{itemize}
Với người tham gia giao thông, các câu hỏi rất cụ thể: nên đi tuyến nào khi xuất phát lúc 7 giờ 45? Khi thấy phía trước
ùn tắc có nên rẽ sang đường khác? Có nên đợi thêm nửa tiếng để tránh cao điểm hay không? Khi rất nhiều người cùng dùng
một ứng dụng chỉ đường, chỉ dẫn ``nhanh nhất'' có còn nhanh nhất?

Song song với đó, các phương pháp học máy --- đặc biệt là học tăng cường (Reinforcement Learning, RL) và mạng nơ-ron
đồ thị (Graph Neural Network, GNN) --- được kỳ vọng nhiều cho định tuyến. RL hứa hẹn học được chính sách ra quyết định
trực tiếp từ trải nghiệm mà không cần mô hình tường minh~\cite{sutton2018}; GNN có khả năng khai thác cấu trúc đồ thị
của mạng đường để dự báo vận tốc~\cite{li2018dcrnn,yu2018stgcn}. Tuy nhiên các phương pháp cổ điển như TD-A* lại có
bảo đảm tối ưu (dưới các giả thiết hợp lý) và chạy rất nhanh. Câu hỏi then chốt vì vậy không phải ``RL có dùng được
không'' mà là \emph{trong những điều kiện nào thì phương pháp học có lợi hơn thuật toán cổ điển, và lợi bao nhiêu}.
Để trả lời công bằng, mọi phương pháp phải được đánh giá trên cùng một mạng đường thật, cùng một tập truy vấn kiểm
tra, cùng một môi trường mô phỏng và cùng một bộ chỉ số --- đó là động lực chính của đề tài.

\subsection{Mục tiêu của đề tài}\label{sec:intro:goal}

\textbf{Mục tiêu chính.} Xây dựng một khung thực nghiệm hoàn chỉnh, chạy được trên máy tính cá nhân chỉ có CPU và
\emph{tái lập được} (cố định seed, toàn bộ tham số nằm trong một tệp cấu hình), đi từ dữ liệu bản đồ mở
OpenStreetMap~\cite{boeing2017} tới báo cáo; trong đó mô phỏng vận tốc theo 24 khung giờ cho mạng đường Hà Đông, cài
đặt và so sánh công bằng các họ phương pháp tìm đường phụ thuộc thời gian, tái lập kế hoạch, dự báo ngắn hạn bằng GNN,
định tuyến theo rủi ro và học tăng cường; đồng thời cung cấp giao diện minh hoạ trực quan.

\textbf{Mục tiêu cụ thể.}
\begin{enumerate}[label=(\arabic*)]
    \item Thu thập và chuẩn hoá mạng đường lái xe quận Hà Đông từ OpenStreetMap, xử lý các tình huống thực tế (ranh giới
    hành chính thay đổi, máy chủ dữ liệu không phản hồi, thiếu thuộc tính tốc độ).
    \item Xây dựng mô hình vận tốc 24 khung giờ có cấu trúc, bảo đảm tính FIFO (đi trước thì đến trước), kèm mô hình
    nhiễu và sự cố cho ``môi trường thực''.
    \item Cài đặt hiệu quả các thuật toán: Dijkstra, A*, TD-Dijkstra, TD-A* với nhiều heuristic, phiên bản biên dịch JIT,
    tái lập kế hoạch (đầy đủ, lười, D* Lite), dự báo bằng mạng tích chập đồ thị và định tuyến theo phân vị.
    \item Thiết kế môi trường học tăng cường chuẩn Gymnasium và huấn luyện nhiều biến thể RL (quan sát cục bộ, quan sát
    cây lân cận với GNN, bộ nhớ LSTM, học bắt chước, định hình thưởng, kết hợp RL với tìm kiếm).
    \item Đánh giá trên tập kiểm tra cố định với kiểm định thống kê; phân tích theo giờ, mức nhiễu, sự cố, khoảng cách;
    các thí nghiệm phụ về chọn giờ xuất phát và về hiệu ứng khi nhiều xe cùng được dẫn đường.
\end{enumerate}

\textbf{Câu hỏi nghiên cứu.}
\begin{itemize}
    \item[\textbf{RQ1.}] Tính đến sự phụ thuộc thời gian (TD-A*) cải thiện thời gian đi bao nhiêu so với tuyến ngắn nhất
    hoặc tuyến theo vận tốc tự do? Lợi ích thay đổi như thế nào theo giờ trong ngày?
    \item[\textbf{RQ2.}] Khi dự báo sai (nhiễu ngày, sự cố), tái lập kế hoạch dựa trên quan sát cục bộ, dự báo GNN từ cảm
    biến, hay định tuyến theo phân vị 90\% giúp được bao nhiêu? Có thể tái lập kế hoạch với chi phí tính toán thấp hơn
    (replan lười, D* Lite) mà không mất chất lượng không?
    \item[\textbf{RQ3.}] RL chọn từng cạnh có cạnh tranh được với tìm kiếm có mô hình không? Các cải tiến (quan sát rộng
    hơn, bộ nhớ, học bắt chước, định hình thưởng) và cách kết hợp RL với tìm kiếm giúp đến đâu?
    \item[\textbf{RQ4.}] Chọn giờ xuất phát có lợi trong điều kiện nào? Điều gì xảy ra khi nhiều xe cùng làm theo một
    chỉ dẫn?
\end{itemize}

\textbf{Giả thuyết.} Từ lý thuyết, chúng tôi kỳ vọng: (H1) TD-A* tốt hơn rõ rệt tuyến tĩnh vì nó tối ưu dưới FIFO;
(H2) tái lập kế hoạch có lợi chủ yếu khi dự báo sai cục bộ và lớn (sự cố); (H3) RL chọn từng cạnh với thông tin cục bộ
khó vượt được tìm kiếm có mô hình, vì nó phải tự học lại cấu trúc đồ thị mà tìm kiếm đã có sẵn; (H4) dẫn đường cá nhân
tối ưu có thể làm hệ thống tệ đi khi tỉ lệ người dùng lớn. Phần~\ref{sec:exp} kiểm chứng từng giả thuyết bằng số liệu.

\subsection{Đóng góp chính của đề tài}\label{sec:intro:contrib}

\begin{enumerate}[label=(\arabic*)]
    \item \textbf{Khung thực nghiệm tái lập được.} Toàn bộ quy trình chạy bằng một lệnh, với hai chế độ (\emph{smoke}
    vài phút để kiểm tra và \emph{full} cho kết quả chính thức); mọi bảng, hình và \emph{mọi con số trong báo cáo này}
    được sinh tự động từ kết quả chạy, không có số liệu gõ tay.
    \item \textbf{Mô hình vận tốc 24 khung giờ cho Hà Đông.} Mô hình có cấu trúc, minh bạch về giả định: profile ùn tắc
    chu kỳ ngày, độ nhạy theo loại đường, các trục nóng, hiệu ứng hướng vào/ra trung tâm, nhiễu tương quan không gian ---
    thời gian; bảo đảm FIFO kể cả khi có sự cố nhờ thủ tục ``FIFO hoá'' có chứng minh.
    \item \textbf{Bộ thuật toán hiệu quả trên cấu trúc CSR.} TD-A* với ba heuristic (khoảng cách chim bay, cận dưới chính
    xác, ALT), phiên bản Numba; ba cơ chế tái lập kế hoạch; dự báo GCN từ 15\% cạnh có cảm biến; định tuyến theo P90.
    \item \textbf{Sáu hướng cải tiến RL trên cùng một môi trường} và một phân tích trung thực về khi nào RL có và không
    có lợi --- bao gồm cả các kết quả âm (RL thua tìm kiếm, định tuyến P90 không có lợi).
    \item \textbf{Hai thí nghiệm phụ} (chọn giờ xuất phát bằng Q-learning; mô phỏng nhiều xe với phản hồi tải BPR) và
    \textbf{ứng dụng minh hoạ} trên nền Streamlit, được đóng gói Docker.
\end{enumerate}

\subsection{Cấu trúc báo cáo}\label{sec:intro:outline}

Phần~\ref{sec:problem} phát biểu bài toán bằng ngôn ngữ tự nhiên và dưới dạng mô hình toán học, kèm ví dụ minh hoạ
và phân tích các khó khăn. Phần~\ref{sec:related} tổng quan các nghiên cứu liên quan, gồm các thuật toán cổ điển và
các hướng tiếp cận hiện đại, kết thúc bằng bảng so sánh và khoảng trống nghiên cứu. Phần~\ref{sec:method} trình bày
chi tiết phương pháp đề xuất: dữ liệu, mô hình mô phỏng, từng thuật toán, môi trường học tăng cường, phân tích độ phức
tạp và tính đúng đắn, cùng sơ đồ hệ thống. Phần~\ref{sec:exp} mô tả thiết lập thực nghiệm và phân tích kết quả.
Phần~\ref{sec:conclusion} tổng kết, nêu hạn chế và hướng phát triển. Phụ lục liệt kê tham số cấu hình và các điểm nóng
ùn tắc được giả định.
